% Compile with XeLaTeX (needs a CJK font for the quoted Korean sentences).
% If XeLaTeX/kotex is unavailable, replace the \korean{...} quotes with their English glosses.
\documentclass[11pt]{article}
\usepackage{kotex}          % Korean support; use xelatex
\usepackage{amsmath, amssymb}
\usepackage[margin=1in]{geometry}
\usepackage{enumitem}
\newcommand{\korean}[1]{\textit{#1}}
\title{A Sentence-by-Sentence Self-Analysis of My Claims\\
on the Tractability of the Bar-Pointer Latent Variable Model}
\author{Jaehoon Ahn}
\date{Prepared in response to Prof.\ Jung's tractability tutorial}
\begin{document}
\maketitle

\section*{0. Purpose and the single most important concession}

Professor, you asked me to quote each of my own sentences and analyze, precisely, why each is
right or wrong---using the framework in your tutorial. I take that seriously, so I have checked
every sentence against Definition~2.1 (tractability $=$ polynomial-time marginal likelihood) and
the five mechanisms of Section~2.5.

Before anything else, the one concession that I believe is the true root of my error:

\begin{quote}
I repeatedly wrote or implied ``\emph{the latent is low-dimensional, therefore exact inference is
possible.}'' \textbf{This causal step is wrong, and you are right to reject it.} Low dimensionality
is only what makes \emph{discretization feasible}; it is a \emph{necessary} condition for one route
to tractability, never a \emph{sufficient} one. The sufficient condition is \emph{structural}:
Markov transitions $+$ factorized emissions (your Mechanism~3), or linear-Gaussian structure (your
Mechanism~4), etc. I collapsed ``low-dimensional'' and ``tractable'' into one idea. They are not one
idea. This is the misconception, stated plainly.
\end{quote}

The rest of this note does two things. First (Sec.~2) it re-derives which structural property each
of my sentences actually depends on. Second, and this is the crux, it shows that the truth value of
my central sentences is \emph{not} a matter of opinion but of \emph{one checkable fact about the
generative model}: whether its transition is Markov and its emission factorized. I state that
criterion here, in your framework, and commit to verifying it in the code line-by-line (Sec.~4). I
do not assert the answer in this document, because I have not yet finished reading every line, and
you are right that this must be established from the code, not asserted.

\section*{1. The framework, in my own words (to show I have absorbed it)}

For a latent-variable model $p_\theta(z,x)=p_\theta(x\mid z)\,p_\theta(z)$, tractability is the
ability to compute the marginal likelihood
\[
  p_\theta(x)=\int p_\theta(x\mid z)\,p_\theta(z)\,dz
\]
in time polynomial in $(K,d,T,N)$ (Def.~2.1). Writing down the numerator pointwise is trivial and
irrelevant; the content is the \emph{normalizer}. Everything downstream (posterior evaluation,
EM's E-step, sampling) follows from having $p_\theta(x)$.

The two mechanisms that matter for our discussion:
\begin{description}[leftmargin=1.4em]
\item[Mechanism 3 (Sec.~2.5.3).] If $z$ is discrete with \emph{Markov} structure and \emph{factorized}
emissions, the $K^{T}$-term sum collapses to $O(TK^2)$ by forward--backward. Discreteness alone is
not enough; the sum collapses \emph{because} of Markov $+$ factorized structure. This is exactly the
distinction you draw between ``expressible'' and ``computable'' (Sec.~2.5.3, the note after Eq.~7).
\item[Mechanism 4 (Sec.~2.5.4).] A $Td$-dimensional linear-Gaussian sequence is tractable by the
Kalman recursions---\emph{high}-dimensional yet tractable. This is your decisive refutation of
``low-dimensional $\Rightarrow$ tractable,'' and I accept it without reservation.
\end{description}

The four-stage progression is the key organizing device, and I now use \emph{your} labels:
\begin{itemize}[leftmargin=1.4em]
\item \textbf{Stage 1:} discrete state, \emph{Markov} transition, \emph{factorized} emission
$\Rightarrow$ Mechanism 3 $\Rightarrow$ tractable (EM/Baum--Welch, Viterbi). This is madmom/Krebs.
\item \textbf{Stage 2--3:} a Transformer \emph{inside the latent prior or the emission} breaks
factorization (Stage 2) or the Markov property of the \emph{latent} chain (Stage 3).
\item \textbf{Stage 4:} continuous $(\varphi,v)$ \emph{on top of} the Stage-3 Transformer prior and
Transformer emission $\Rightarrow$ intractable; VAE machinery is then a technical necessity.
\end{itemize}

\section*{2. Sentence-by-sentence analysis}

For each sentence I give: the verdict, and the structural property it depends on.

\paragraph{S1. \korean{``이 모델은 HMM이 아니다; HMM이라 부른 적도 없다.''}}
(``This model is not an HMM.'') \textbf{Correct as stated.} A continuous-state model is literally not
an HMM (HMMs have finite discrete states). This is a true factual statement and carries no claim
about tractability by itself.

\paragraph{S2. \korean{``$z$가 continuous이므로 정확 추론은 여전히 intractable하다.''}}
\textbf{Correct---and it is a concession in your favor.} With a wrapped/von~Mises transition and a
cosine-bump (nonlinear) emission, the continuous posterior is neither conjugate (Mech.~1 fails) nor
linear-Gaussian (Mech.~4 fails). So exact inference \emph{on the continuous latent} is intractable.
I agree.

\paragraph{S3. \korean{```Viterbi 돌릴 만큼 단순하다'는 과장이었지만 일리는 있다.''}}
\textbf{Correct/calibrated, with one clarification I owe you.} The clarification: for \emph{training}
the relevant object is the forward algorithm (the marginal likelihood, Mech.~3); Viterbi is the
\emph{deployment-time} decoder. My shorthand ``Viterbi'' conflated the two. I was \emph{not}
confusing training with inference conceptually---my intended claim was ``the posterior is tractable,
hence trainable by an EM-type exact-likelihood method rather than by ELBO/VAE,'' which is precisely
your own dichotomy (your message, point~1). But the word ``Viterbi'' was the wrong token for
``forward algorithm,'' and I accept the correction.

\paragraph{S4. \korean{``잠재공간이 저차원적이다.''}}
\textbf{Correct about the per-frame state} $(\varphi,v,m)\in\mathbb{R}^2\times\{2,\dots,7\}$, which is
what makes per-state discretization into $K$ bins feasible. Your \S8.2 nuance (the \emph{sequence}
latent is $3T$-dimensional) is true but does not falsify S4, because no one discretizes the $3T$-dim
sequence space directly---that would be $K^T$. Under Mechanism 3 one discretizes the \emph{per-frame}
state and lets forward--backward handle $T$ in $O(TK^2)$. So the relevant quantity is exactly the
per-frame dimension, and it is low. (But per S0, low dimension alone proves nothing; it only enables
the discretization \emph{route}.)

\paragraph{S5. \korean{``phase/tempo/meter는 이미지 잠재보다 천만 배 단순하다.''}}
\textbf{The factual content is correct} (3 physical dimensions vs.\ hundreds/thousands for image
latents; ``천만 배'' is rhetorical). \textbf{The \emph{implication} I attached to it was wrong}: I
used it to suggest tractability, which is the S0 error. Restated correctly: the low per-frame
dimension makes the \emph{discretization route} available; whether that route yields tractability
still requires the structural condition (Mech.~3).

\paragraph{S6 (the crux). \korean{``지금 설계된 모델은 HMM은 아니지만, 이산화만 하면 HMM으로
축소 가능하다. Image VAE에서는 성립하지 않는다.''}}
\textbf{This sentence is conditionally true, and the condition is a checkable property of the
generative model---not a matter of opinion.} Precisely:
\begin{itemize}[leftmargin=1.4em]
\item \emph{If} the generative model is Stage-1-structured---transition
$p(z_t\mid z_{t-1}, h_t)$ that is Markov in $z$ (where $h_t$ is a fixed function of the
\emph{observed} audio $x$, not of the latent $z$), and emission $p(b_t\mid z_t)$ that is per-frame
factorized---\emph{then} discretizing $(\varphi,v)$ yields a genuine (time-inhomogeneous,
input-conditioned) HMM, and Mechanism 3 applies. S6 is \emph{true}.
\item \emph{If} the generative model is Stage-3/4---a Transformer \emph{prior over the latent
sequence} $z_{1:T}$ (non-Markov in $z$) or a Transformer \emph{emission} attending over $z_{1:T}$
(non-factorized)---\emph{then} your \S8.1 is exactly right: discretizing $(\varphi,v)$ recovers only
the discrete \emph{state space} while leaving the non-Markov/non-factorized structure intact, so it
does \emph{not} give an HMM. S6 is \emph{false}.
\end{itemize}
The image-VAE half of S6 is unconditionally correct for a separate reason: an image latent of
hundreds of dimensions cannot be gridded at all ($K^{\text{hundreds}}$), so the discretization route
is unavailable regardless of structure.

\medskip
\noindent\textbf{Therefore the entire disagreement reduces to one question about my code:}
\emph{is the transition Markov in $z$ (conditioned on audio-derived context), and is the emission
per-frame factorized?} I believe the \emph{design intent} is the Stage-1 structure---the Transformer
in the model processes the \emph{audio} $x$ into a context sequence $h_{1:T}$, and the latent
transition consumes $h_t$ and $z_{t-1}$ only---which would make the model an input-conditioned
Stage-1 model rather than a Stage-4 model. But I have not yet audited every line, and you are right
that this must be established from the implementation, not from my description. Section~4 is my
verification plan.

\paragraph{S7. \korean{``exact inference가 가능하면 VAE를 쓸 이유가 없다; VAE는 $p(z|x)$가
intractable일 때의 기법이다.''}}
\textbf{The conclusion is correct and is your own point 4}: the VAE (ELBO $+$ amortized encoder)
exists to approximate an \emph{intractable} posterior; where exact inference is available, EM-type
exact-likelihood training is the appropriate tool (GMM is the canonical example you cite). \textbf{The
one weak link is the antecedent} ``because the latent is low-dimensional, exact inference is
possible''---the S0 error again. Corrected: ``\emph{if} the model is Stage-1-structured (Mech.~3),
exact inference is possible, and then a VAE is unnecessary.''

\paragraph{S8. \korean{``이미지 VAE의 잠재공간은 학습 후에야 정의되지만, bar-pointer의 잠재는
이미 well-defined 물리량이다.''}}
\textbf{True as an observation, and I accept your correction about its scope.} You are right (\S2.6,
your point 2) that physical meaning is \emph{orthogonal to tractability}: an image VAE and a
neurally-parameterized bar-pointer VAE can both be intractable. So S8 is \emph{not} a tractability
argument, and I should not have deployed it as one. Its correct role is narrower and separate: it
explains why the \emph{encoder} is less essential---in an image VAE the latent semantics must be
\emph{discovered} by training, whereas here the latent targets are pre-specified, so there is nothing
for a recognition network to ``invent.'' That is a statement about the encoder's necessity, not
about the marginal likelihood.

\section*{3. Consolidated: exactly where I was wrong}

\begin{enumerate}[leftmargin=1.6em]
\item \textbf{Dimension $\Rightarrow$ tractability (S5, S7 antecedent, and the spirit of S4/S6).}
The genuine misconception. Tractability needs \emph{structure} (Mech.~3 or 4), not smallness. Your
Kalman example (Mech.~4: 1000-dim, tractable) is the clean disproof and I accept it fully.
\item \textbf{``Viterbi'' as shorthand for the training objective (S3).} The training object is the
forward algorithm / marginal likelihood; Viterbi is deployment decoding. Sloppy tokenization on my
part.
\item \textbf{Using physical meaning as a tractability argument (S8).} It is orthogonal to
tractability; it bears only on the encoder's necessity.
\end{enumerate}

\section*{4. What is \emph{not} settled by opinion, and my verification plan}

Sentences S1--S5, S7, S8 are settled above on their own terms. The crux, S6, is settled \emph{only}
by a property of the code, and I will not claim it until I have verified it. The criterion, stated in
your framework, is a two-part checklist:
\begin{enumerate}[leftmargin=1.6em]
\item \textbf{Markov test.} In the generative prior, does $z_t$ depend on $z_{1:t-1}$, or only on
$z_{t-1}$ and an audio-derived context $h_t=f(x)$? Only the latter is Markov-in-$z$ (Stage 1); the
former is Stage 3. (Conditioning the transition \emph{parameters} on all of $x$ via a bidirectional
$h$ does not break Markov-in-$z$; it is an input-conditioned HMM.)
\item \textbf{Factorization test.} Does the emission $p(b_t\mid\cdot)$ read only $z_t$, or does it
attend over $z_{1:T}$? Only the former is factorized (Stage 1).
\end{enumerate}
I will audit \texttt{rollout\_prior} and the emission module line-by-line against this checklist and
report which stage the implementation actually realizes. If it is Stage 1, S6 is true and Mechanism 3
applies; if a Transformer sits inside the latent prior/emission, S6 is false and your \S8.1 stands.
As an \emph{independent} empirical check of the same fact, I will run the forward algorithm
(Mechanism 3's own algorithm) on the model: if it computes a well-defined marginal likelihood and
trains the model, the generative structure is necessarily Markov$+$factorized, because forward
algorithm is not merely inefficient but \emph{ill-defined} on a Stage-4 model. (I will present those
numbers separately, after the code audit, so that the argument here rests on structure, not on
results.)

\section*{5. Agreement on the path forward}

Your proposed plan---formulate the bar-pointer model as a Stage-1 HMM, train its parameters by EM,
publish that, and \emph{then} extend toward the VAE---is, I think, exactly right, for the reason this
note makes precise: Stage~1 is the tractable regime (Mech.~3), and it is where exact-likelihood
training lives. I would like to build it in that order. Implementing Stage~1 with EM will also be the
cleanest way for me to internalize, at the level of code, the distinction between ``discrete state''
and ``Markov$+$factorized structure'' that I blurred---which is the concept I most need to solidify.

\medskip
\noindent I am aware these are not easy ideas and that my probability/statistics vocabulary has been
imprecise. I am treating that as a gap to close, not a verdict. Thank you for writing the tutorial;
using its framework to check my own sentences is the exercise that finally located my error precisely
(Sec.~3), which no amount of back-and-forth in text had done.

\end{document}
